Since $c^2\ge10$, we have $c>3$. If $a=0$, then $c$ is even and $(c-3)(c+3)=3^b$. The odd factors are powers of $3$ differing by $6$, necessarily $3,9$. This gives $(a,b,c)=(0,3,6)$.

For $a\ge1$, $c$ is odd. Modulo $8$, $a=1$ or $2$ is impossible, so $a\ge3$. Put $u=\dfrac{c-3}{2}$ and $v=\dfrac{c+3}{2}$. Then $v-u=3$ and $uv=2^{a-2}3^b$.
If $b=0$, the opposite-parity factors are powers of two, hence $u=1,v=4$, giving $(4,0,5)$. If $b=1$, divisibility by $3$ of one factor forces it for both, impossible. For $b\ge2$, put $u=3r,v=3s$. Then $s-r=1$ and $rs=2^{a-2}3^{b-2}$.
When $b=2$, consecutive powers-of-two factors must be $1,2$, giving $(3,2,9)$.
When $b>2$, coprimality forces $r,s$ to be the powers $2^{a-2},3^{b-2}$ in some order. The positive-exponent solutions of $|2^m-3^n|=1$ are $(m,n)=(2,1),(1,1),(3,2)$. To verify this elementary fact, the equation $2^m-3^n=1$ forces $m\le2$ modulo $8$; in the reverse equation, $m=1$ gives $(1,1)$, while $m\ge3$ forces even $n$ and $(3^{\dfrac{n}{2}}-1)(3^{\dfrac{n}{2}}+1)=2^m$, whose factors must be $2,4$.
These three pairs give $(4,3,21),(3,3,15),(5,4,51)$.
All six triples check directly, so the answer is $\boxed6$.
